\chapter{Theoretischer Hintergrund}
\label{ch:hauptteil}

\section{Grundlagen der Sozialen Arbeit}

Die Soziale Arbeit versteht sich als Handlungswissenschaft, die sowohl theoretische als auch praktische Dimensionen umfasst \parencite[vgl.][87]{thiersch2014}. Nach \textcite[15]{staub2018} ist die Lebensweltorientierung ein zentrales Paradigma.

\subsection{Lebensweltorientierung}

Der Ansatz der Lebensweltorientierten Sozialen Arbeit wurde maßgeblich von Hans Thiersch geprägt. Er betont die Notwendigkeit, \enquote{die Alltagserfahrungen der Adressaten ernst zu nehmen und als Ausgangspunkt professionellen Handelns zu verstehen} \parencite[23]{thiersch2014}.

\subsection{Professionalisierungsdebatte}

Die Diskussion um die Professionalisierung der Sozialen Arbeit wird seit Jahrzehnten geführt. Tabelle~\ref{tab:vergleich} gibt einen Überblick über zentrale Positionen.

\begin{table}[htbp]
    \centering
    \caption{Vergleich zentraler Positionen zur Professionalisierung}
    \label{tab:vergleich}
    \begin{tabularx}{\textwidth}{lXX}
        \toprule
        \textbf{Autor/in} & \textbf{Position} & \textbf{Kernaussage} \\
        \midrule
        Müller (2020) & Kompetenzorientiert & Professionalisierung durch Standards \\
        Staub-Bernasconi (2018) & Menschenrechtsbasiert & Tripel-Mandat als Grundlage \\
        Thiersch (2014) & Lebensweltorientiert & Alltagsnähe als Qualitätsmerkmal \\
        \bottomrule
    \end{tabularx}
    \source{Eigene Darstellung.}
\end{table}

\section{Methodische Grundlagen}

Abbildung~\ref{fig:prozess} zeigt den typischen Forschungsprozess in der qualitativen Sozialforschung.

\begin{figure}[htbp]
    \centering
    \begin{tikzpicture}[
        box/.style={rectangle, draw, rounded corners, minimum width=3.5cm, minimum height=1cm, align=center},
        arr/.style={-{Stealth}, thick}
    ]
        \node[box] (a) {Fragestellung};
        \node[box, right=1.5cm of a] (b) {Datenerhebung};
        \node[box, right=1.5cm of b] (c) {Auswertung};
        \node[box, below=1cm of c] (d) {Interpretation};
        \draw[arr] (a) -- (b);
        \draw[arr] (b) -- (c);
        \draw[arr] (c) -- (d);
    \end{tikzpicture}
    \caption{Forschungsprozess in der qualitativen Sozialforschung}
    \label{fig:prozess}
    \source{Eigene Darstellung in Anlehnung an \textcite[58]{schmidt2019}.}
\end{figure}

\section{Zwischenfazit}

Die dargestellten theoretischen Grundlagen zeigen, dass die Soziale Arbeit über ein breites methodisches Repertoire verfügt. Für die vorliegende Untersuchung ist insbesondere der lebensweltorientierte Ansatz relevant, da er die Perspektive der Adressatinnen und Adressaten in den Mittelpunkt stellt.
